% tgp_time_glossary.tex -- Glossary of time concepts and EEP analysis
% Self-contained section for \input{} into main manuscript.
% Synchronized with sek08c, sek08 (rem:substrate-time), sek04 (ax:c).

\subsection{Glossary: substrate time, proper time, and clocks}%
\label{ssec:time-glossary}

TGP introduces a non-standard time structure.  To prevent confusion,
we define all temporal quantities explicitly.

\begin{center}\small
\begin{tabular}{@{}lp{4.5cm}p{6.5cm}@{}}
\toprule
\textbf{Symbol} & \textbf{Name} & \textbf{Definition and role} \\
\midrule
$dt_{\rm sub}$
  & Substrate tick
  & Global, $\Phi$-independent increment of the discrete substrate clock.
    Not directly measurable; serves as the ``absolute time'' of the lattice
    $\Gamma$.  All observers count the same number of substrate ticks
    between two events. \\[6pt]
$d\tau$
  & Proper time
  & Time measured by a physical clock (e.g.\ Cs hyperfine transition).
    Depends on local $\Phi$:
    \begin{equation*}
      d\tau = \frac{c_0}{\sqrt{\psi}}\,dt_{\rm sub},
      \qquad \psi \equiv \Phi/\Phi_0.
    \end{equation*}
    In regions with $\Phi > \Phi_0$ (``denser space''): $d\tau < c_0\,dt_{\rm sub}$
    --- clocks run slower. \\[6pt]
$c_{\rm lok}$
  & Local speed of light
  & Measured by a \emph{local} observer using their proper time and
    proper length:
    $c_{\rm lok} = c_0/\sqrt{\psi}$.
    This is \textbf{not} the same as the coordinate speed
    $c_{\rm coord} = c_0/\psi$ (which is frame-dependent). \\[6pt]
$c_0$
  & Reference speed
  & Speed of light at $\Phi = \Phi_0$ (reference vacuum).
    Identified with the SI value $c = 299\,792\,458$\,m/s.
    In TGP, $c_0$ is a \emph{property of the reference state},
    not a universal constant. \\[6pt]
$\nu_{\rm Cs}(\Phi)$
  & Cs clock frequency
  & Hyperfine transition frequency of $^{133}$Cs (SI second definition).
    In TGP: $\nu_{\rm Cs}(\Phi) = \nu_{\rm Cs,0}\,\sqrt{\Phi_0/\Phi}$.
    Two clocks at different $\Phi$ tick at different rates
    --- this is the gravitational redshift, here derived from substrate
    budget rather than postulated. \\[6pt]
$T_{\rm period}$
  & Orbital/dynamical period
  & Period of a macroscopic orbit (e.g.\ planet).
    Measured in substrate ticks, it is $\Phi$-independent.
    Measured in proper time, it depends on $\Phi$. \\
\bottomrule
\end{tabular}
\end{center}

\paragraph{Einstein Equivalence Principle (EEP).}
TGP maintains a \emph{modified} form of EEP:
\begin{enumerate}[nosep]
  \item \textbf{Weak Equivalence Principle (WEP):}
        \textsc{Satisfied.}  In TGP the inertial mass equals the
        gravitational source strength ($m_{\rm inert} = C_i = m_{\rm grav}$);
        all bodies fall identically in a given $\nabla\Phi$.
  \item \textbf{Local Lorentz Invariance (LLI):}
        \textsc{Satisfied locally.}  In a sufficiently small region,
        $\Phi \approx \text{const}$ and the metric reduces to Minkowski
        with $c = c_{\rm lok}$.  No preferred direction arises.
  \item \textbf{Local Position Invariance (LPI):}
        \textsc{Modified.}  Non-gravitational experiments
        (e.g.\ fine structure constant $\alpha$, mass ratios)
        are $\Phi$-independent in TGP --- they depend on
        dimensionless combinations only.  But $c$, $\hbar$, $G$
        individually depend on $\Phi$.
        The \emph{observable} consequence: gravitational redshift
        $z_{\rm grav} = \sqrt{\Phi_{\rm emit}/\Phi_{\rm obs}} - 1$
        follows the \textbf{same} formula as GR (to leading order
        in $\delta\Phi/\Phi_0$).
\end{enumerate}

\paragraph{Operational protocol.}
To compare TGP predictions with clock experiments
(GPS, Gravity Probe A, ACES, MICROSCOPE):
\begin{enumerate}[nosep]
  \item Express all observables as \textbf{dimensionless frequency ratios}:
        $\nu_A/\nu_B$ between clocks A and B.
  \item In TGP: $\nu_A/\nu_B = \sqrt{\Phi_B/\Phi_A} = \sqrt{\psi_B/\psi_A}$.
  \item In GR: $\nu_A/\nu_B = \sqrt{g_{tt}^{(B)}/g_{tt}^{(A)}}
        \approx 1 + (U_A - U_B)/c^2$.
  \item To leading order in $\delta\Phi/\Phi_0$: \textbf{identical}.
        Deviations appear at $O(\delta\Phi^2/\Phi_0^2)$, i.e.\
        $O(U^2/c^4) \sim 10^{-18}$ on Earth --- currently below
        experimental sensitivity.
\end{enumerate}

\paragraph{Relation to SI.}
In the 2019 SI, the values of $c$ and $h$ are \emph{exact by definition}:
$c = 299\,792\,458$\,m/s, $h = 6.626\,070\,15 \times 10^{-34}$\,J\,s.
TGP does \emph{not} contradict this: the SI definitions fix the
\emph{numerical value} of $c$ in the unit system, regardless of whether
$c$ is a universal constant or a field.  The physical content of TGP
is encoded in \textbf{dimensionless observables} (frequency ratios,
fine structure constant, mass ratios) that are measurement-convention-independent.
